\documentclass[10pt,a4paper]{amsart}
\usepackage[margin=19mm]{geometry}
\usepackage{amsmath,amssymb,mathtools}
\usepackage[scr=rsfs,cal=euler]{mathalfa}
\usepackage{xfrac}
\usepackage[unicode,hidelinks,bookmarks=false]{hyperref}
\newcommand{\ie}{i.e.\ }
\newcommand{\cf}{cf.\ }
\newcommand{\resp}{respectively\ }
\newcommand{\Subsection}{\subsection}
\providecommand{\nobreakdash}{\nobreak\hbox{-}}
\setlength{\emergencystretch}{2em}
\multlinegap0pt
\usepackage{bm}
\usepackage{bbm}
\usepackage{dsfont}
\usepackage{xspace}
\usepackage{cancel}
\usepackage{mathtools}
\usepackage{amsthm}
\makeatletter
\@ifundefined{theorem}{\newtheorem{theorem}{Theorem}}{}
\@ifundefined{lemma}{\newtheorem{lemma}[theorem]{Lemma}}{}
\makeatother
\makeatletter
\expandafter\ifx\csname subproof\endcsname\relax
\newenvironment{subproof}[1][\proofname]{%
  \renewcommand{\qedsymbol}{$\spadesuit$}%
  \begin{proof}[#1]%
}{%
  \end{proof}%
}
\fi
\makeatother
\title[How many colours to connect cliques?]{How many colours to connect cliques?}
\author{A. Manas}
\date{Source: arXiv:2608.06418v1 (CC BY 4.0)}
\begin{document}
\maketitle

\noindent\emph{Source and licence.} How many colours to connect cliques?. Version arXiv:2608.06418v1 (CC BY 4.0), distributed under the Creative Commons Attribution 4.0 International licence, as recorded in the arXiv licence metadata for this exact version and shown on the abstract page https://arxiv.org/abs/2608.06418v1. Full rights notice: licenses/arxiv-2608.06418-CC-BY-4.0.txt.\par\smallskip

\noindent\textit{Editorial scope.} Counted unit: the Lemma whose source label is \texttt{lem:3:2} in the article (source0.tex lines 232--235, source label \texttt{lem:3:2}), delivered together with the article's own complete original proof of it (source0.tex lines 236--255, one \texttt{proof} environment of 20 lines). No mathematics is written, repaired or re-derived: statement and proof are the article's own text line for line. Required context retained uncounted: none. Every definition, display and label of those lines is retained unchanged. Omitted from this package and not redistributed: 1--231, 256--383. Nothing inside a retained excerpt is omitted or altered: every retained body is delivered exactly as the article's lines are written, so no omission receipt of any kind accompanies this package. Citations: the retained excerpts carry no citation command at all, so no bibliography entry of the article is transcribed and no reference is redistributed. Reference adaptation: none was needed and none was made; every reference inside the delivered unit resolves inside it, so no unresolved reference or citation is produced. Printed slips kept literal: no printed word, symbol, hypothesis or number of the counted statement or of its proof was altered. Layout and class: the article's own class is \texttt{article} with packages including amsmath, amsthm, nicefrac, amsfonts, amstext, graphicx, textcomp, caption, enumitem, mathrsfs, amssymb, xcolor, esvect, algorithm2e; the wrapper is the lane's pinned \texttt{amsart} editorial wrapper at 10pt on A4, which restates the theorem-like environments the retained text needs and adds the article's own macro definitions that the retained text needs (none), with extra wrapper packages (bm, bbm, dsfont, xspace, cancel, mathtools). The delivered numbering is the unit's own. Assets: no retained span contains a figure environment, a graphics-inclusion command, a PDF-page inclusion, a table or a TikZ picture; no image file is copied into this package, the package declares no asset dependency and no image file is redistributed with it. Licence: version arXiv:2608.06418v1 of this article is distributed under the Creative Commons Attribution 4.0 International licence, as recorded in the arXiv licence metadata for this exact version and shown on the abstract page https://arxiv.org/abs/2608.06418v1. A full rights notice is delivered at licenses/arxiv-2608.06418-CC-BY-4.0.txt. Characters: every retained excerpt is pure ASCII, so the delivered engine, XeLaTeX with its default Latin Modern text font, reports no missing character.
\begin{lemma}
\label{lem:3:2}
Let $a, b$ be positive integers. Then $B\left(a\binom{2b+1}{b+1}, \binom{a+1}{2}(2b+1)\right) \geq \frac{1}{2}a(b+1)$.
\end{lemma}

\begin{proof}
Let $\mathcal F$ denote the set of all $(b+1)$-element subsets of $[2b+1]$, and let $\mathcal G$ denote the set of all $2$-element subsets of $\{0\} \cup[a]$.  Our vertex set will be given by $[a] \times \mathcal F$, while our colours will come from $\mathcal G \times [2b+1]$. \\

\noindent
The colouring $\chi : \binom{[a] \times \mathcal F}{2} \mapsto \mathcal G \times [2b+1]$ is described by:
\begin{itemize}
    \item $\chi((i, S_1), (i, S_2)) = (\{0,i\}, \min(S_1 \cap S_2))$ for all $i \in [a]$ and $S_1 \ne S_2 \in \mathcal F$.
    \item $\chi((i, S_1), (j, S_2)) = (\{i,j\}, \min(S_1 \cap S_2))$ for all $i \ne j \in [a]$ and $S_1, S_2 \in \mathcal F$.
\end{itemize}
We note that $\chi$ is well-defined since any two elements of $\mathcal F$ intersect. Let us also remark that the choice of ``$\min$'' for the second coordinate is arbitrary, the argument that follows only requires that this coordinate be contained in $S_1 \cap S_2$. \\

\noindent
Let $C \subseteq \mathcal{G} \times [2b+1]$ be a set of colours such that each vertex has an incident edge with colour in $C$, that is, no vertex is left isolated. Denote by $X$ the set of pairs $(x, c) \in [a] \times [2b+1]$ for which there exists $x'$ so that $(\{x, x'\}, c) \in C$. Note that each element $(\{x, y\}, c)$ gives rise to at most two such pairs in $X$, namely, $(x, c)$ and $(y, c)$ provided both $x, y \ne 0$. Hence, $|X| \leq 2|C|$. \\

\noindent
On the other hand, we claim that for each $x \in [a]$ there must be at least $b+1$ elements $c$ of $[2b+1]$ with $(x,c) \in X$. If this were not the case, then we could find a set $S$ with at least $(2b+1)-b = b+1$ elements so that no element in $X$ is of the form $(x, c)$ with $c \in S$. However, this would imply that $(x, S)$ is isolated, a contradiction. \\

\noindent
It follows that $2|C| \geq |X| \geq a \cdot (b+1)$, and thus to ``cover'' all vertices one must use at least $\frac{1}{2}a(b+1)$ colours, proving the lemma.
\end{proof}

\end{document}
